\centerline{\bf \Huge I Тренировочная Олимпиада}

\begin{flushright}
        \scriptsize{\textbf{Именно математика даёт надёжнейшие правила: тому кто им следует — тому не опасен обман чувств.\\
         Леонард Эйлер}}
        
\end{flushright}

\problem{\textbf{1}} Дан числовой ребус:

С + У + М + М + А = 11

Чему может быть равна буква М, если среди букв нет нуля?

\problem{\textbf{2}} Кошка Расула Владимировича родила котят! Известно, что два самых
легких весят в сумме 80 г., четыре самых тяжелых – 200 г., а
суммарный вес всех котят равен 500г. Cколько котят родила кошка?

\problem{\textbf{3}} АА загадал число $X$. Известно, что $X$ делится на 4 и оно меньше 2025. Сколько различных чисел $X$ мог загадать АА, если он не использовал цифры 1, 3, 6, 7, 0?

\problem{\textbf{4}} Между городами A и B ездят автобусы с одинаковой постоянной скоростью. Автобус,
выехавший из A в полдень, и автобус, выехавший из B в 15.00, встретились на расстоянии 500км от A. Автобус, выехавший из A в 14.00, и автобус, выехавший из B в 11.00, встретились на
расстоянии 300 км от A. На каком расстоянии от A встретятся автобусы, выехавшие из A и из B
в 13.00?

\problem{\textbf{5}} Квадратное поле размером 64 квадратных метра разделили на квадратные и прямоугольные
участки размером соответственно 2 × 2 и 4 × 1 метра. Оказалось, что общая длина границы
составила 54 метра (без учета внешних границ поля). Сколько получилось квадратных и
прямоугольных участков?

\problem{\textbf{6}} Двое кладоискателей нашли большой алмаз, разбитый на несколько кусков. Докажите, что один
из его кусков можно разрезать на два так, чтобы весь алмаз можно было разделить на две части
равного веса.